\section{Some Basics of Representation of Lie algebra}

\begin{definition}
    Let \(L\) be a Lie algebra over a field \(F\), and let \(V\) be a vector space over \(F\).
    An \emph{\(L\)-module structure} on \(V\) is a bilinear map
    \[
        L\times V \to V,\qquad (x,v)\mapsto x\cdot v,
    \]
    such that
    \[
        [x,y]\cdot v = x\cdot (y\cdot v)-y\cdot (x\cdot v)
    \]
    for all \(x,y\in L\) and \(v\in V\).
\end{definition}

\begin{definition}
    Let \(V\) and \(W\) be \(L\)-modules. The \emph{tensor product module}
    \(V\otimes W\) is the vector space \(V\otimes_F W\) equipped with the action
    \[
        x\cdot (v\otimes w)=(x\cdot v)\otimes w+v\otimes (x\cdot w),
        \qquad x\in L,\ v\in V,\ w\in W,
    \]
    extended linearly to all of \(V\otimes W\).
\end{definition}

\begin{proof}
    We verify that this defines an \(L\)-module structure. Since the formula is bilinear in
    \(x\), \(v\), and \(w\), it extends uniquely to a linear action on \(V\otimes W\).
    It remains to check the bracket identity.

    For \(x,y\in L\), \(v\in V\), and \(w\in W\), we have
    \begin{align*}
        x\cdot \bigl(y\cdot (v\otimes w)\bigr)
         & =x\cdot \bigl((y\cdot v)\otimes w+v\otimes (y\cdot w)\bigr) \\
         & =(x\cdot (y\cdot v))\otimes w+(y\cdot v)\otimes (x\cdot w)
        +(x\cdot v)\otimes (y\cdot w)+v\otimes (x\cdot (y\cdot w)),
    \end{align*}
    and similarly
    \begin{align*}
        y\cdot \bigl(x\cdot (v\otimes w)\bigr)
         & =(y\cdot (x\cdot v))\otimes w+(x\cdot v)\otimes (y\cdot w)
        +(y\cdot v)\otimes (x\cdot w)+v\otimes (y\cdot (x\cdot w)).
    \end{align*}
    Subtracting, we obtain
    \begin{align*}
        x\cdot \bigl(y\cdot (v\otimes w)\bigr)-y\cdot \bigl(x\cdot (v\otimes w)\bigr)
         & =\bigl([x,y]\cdot v\bigr)\otimes w+v\otimes \bigl([x,y]\cdot w\bigr) \\
         & =[x,y]\cdot (v\otimes w).
    \end{align*}
    Thus \(V\otimes W\) is an \(L\)-module.
\end{proof}

\begin{definition}
    Let \(V\) and \(W\) be \(L\)-modules. The vector space
    \[
        \operatorname{Hom}_F(V,W)
    \]
    becomes an \(L\)-module by defining
    \[
        (x\cdot f)(v)=x\cdot f(v)-f(x\cdot v),
        \qquad x\in L,\ f\in \operatorname{Hom}_F(V,W),\ v\in V.
    \]
    This is called the \emph{Hom-module}.
\end{definition}

\begin{proof}
    We first note that for fixed \(x\in L\) and \(f\in \operatorname{Hom}_F(V,W)\),
    the map \(x\cdot f\) is again \(F\)-linear in \(v\), so \(x\cdot f\in \operatorname{Hom}_F(V,W)\).

    Now let \(x,y\in L\), \(f\in \operatorname{Hom}_F(V,W)\), and \(v\in V\). Then
    \begin{align*}
        \bigl(x\cdot (y\cdot f)\bigr)(v)
         & =x\cdot \bigl((y\cdot f)(v)\bigr)-(y\cdot f)(x\cdot v)                            \\
         & =x\cdot \bigl(y\cdot f(v)-f(y\cdot v)\bigr)
        -\bigl(y\cdot f(x\cdot v)-f(y\cdot (x\cdot v))\bigr)                                 \\
         & =x\cdot (y\cdot f(v))-x\cdot f(y\cdot v)-y\cdot f(x\cdot v)+f(y\cdot (x\cdot v)),
    \end{align*}
    and similarly
    \begin{align*}
        \bigl(y\cdot (x\cdot f)\bigr)(v)
         & =y\cdot (x\cdot f(v))-y\cdot f(x\cdot v)-x\cdot f(y\cdot v)+f(x\cdot (y\cdot v)).
    \end{align*}
    Subtracting gives
    \begin{align*}
        \bigl(x\cdot (y\cdot f)-y\cdot (x\cdot f)\bigr)(v)
         & =[x,y]\cdot f(v)-f([x,y]\cdot v) \\
         & =\bigl([x,y]\cdot f\bigr)(v).
    \end{align*}
    Hence
    \[
        [x,y]\cdot f=x\cdot (y\cdot f)-y\cdot (x\cdot f),
    \]
    so \(\operatorname{Hom}_F(V,W)\) is an \(L\)-module.
\end{proof}

\begin{definition}
    Let \(V\) be an \(L\)-module. The dual space
    \[
        V^*=\operatorname{Hom}_F(V,F)
    \]
    is called the \emph{dual module} of \(V\), where \(F\) is regarded as the trivial
    \(L\)-module. Its action is given by
    \[
        (x\cdot \varphi)(v)=-\,\varphi(x\cdot v),
        \qquad x\in L,\ \varphi\in V^*,\ v\in V.
    \]
\end{definition}

\begin{proof}
    This is the special case \(W=F\) of the Hom-module construction, where the action of
    \(L\) on \(F\) is trivial. Indeed, if \(\varphi\in \operatorname{Hom}_F(V,F)\), then
    \[
        (x\cdot \varphi)(v)=x\cdot \varphi(v)-\varphi(x\cdot v)=0-\varphi(x\cdot v)
        =-\varphi(x\cdot v).
    \]
    Thus \(V^*\) is an \(L\)-module.
\end{proof}

\begin{remark}
    The subspace of \(L\)-invariants in \(\operatorname{Hom}_F(V,W)\) is precisely
    \[
        \operatorname{Hom}_L(V,W)
        =
        \{f\in \operatorname{Hom}_F(V,W)\mid x\cdot f=0 \text{ for all } x\in L\}.
    \]
    Indeed, the condition \(x\cdot f=0\) means
    \[
        x\cdot f(v)=f(x\cdot v)
        \qquad \text{for all } x\in L,\ v\in V,
    \]
    which is exactly the condition that \(f\) be an \(L\)-module homomorphism.
\end{remark}

\begin{proposition}
    Assume that \(V\) is finite-dimensional. Then there is a natural vector space isomorphism
    \[
        V^*\otimes W \cong \operatorname{Hom}_F(V,W),
    \]
    given by
    \[
        \varphi\otimes w \longmapsto \bigl(v\mapsto \varphi(v)w\bigr).
    \]
    This is in fact an isomorphism of \(L\)-modules.
\end{proposition}

\begin{proof}
    Let
    \[
        \Phi:V^*\otimes W \to \operatorname{Hom}_F(V,W)
    \]
    be defined by
    \[
        \Phi(\varphi\otimes w)(v)=\varphi(v)w.
    \]
    It is the usual canonical isomorphism of vector spaces when \(V\) is finite-dimensional.
    We check that it is \(L\)-equivariant.

    For \(x\in L\), \(\varphi\in V^*\), \(w\in W\), and \(v\in V\), we have
    \begin{align*}
        \Phi\bigl(x\cdot (\varphi\otimes w)\bigr)(v)
         & =\Phi\bigl((x\cdot \varphi)\otimes w+\varphi\otimes (x\cdot w)\bigr)(v) \\
         & =(x\cdot \varphi)(v)\,w+\varphi(v)\,(x\cdot w)                          \\
         & =-\varphi(x\cdot v)\,w+\varphi(v)\,(x\cdot w).
    \end{align*}
    On the other hand,
    \begin{align*}
        \bigl(x\cdot \Phi(\varphi\otimes w)\bigr)(v)
         & =x\cdot \bigl(\Phi(\varphi\otimes w)(v)\bigr)-\Phi(\varphi\otimes w)(x\cdot v) \\
         & =x\cdot (\varphi(v)w)-\varphi(x\cdot v)w                                       \\
         & =\varphi(v)(x\cdot w)-\varphi(x\cdot v)w.
    \end{align*}
    Thus
    \[
        \Phi\bigl(x\cdot (\varphi\otimes w)\bigr)=x\cdot \Phi(\varphi\otimes w),
    \]
    so \(\Phi\) is an \(L\)-module isomorphism.
\end{proof}

\begin{remark}
    Under the identification
    \[
        \operatorname{End}_F(V)=\operatorname{Hom}_F(V,V),
    \]
    the \(L\)-action is
    \[
        (x\cdot T)(v)=x\cdot T(v)-T(x\cdot v),
        \qquad x\in L,\ T\in \operatorname{End}_F(V),\ v\in V.
    \]
    Thus \(L\) acts on \(\operatorname{End}_F(V)\) by commutator.
\end{remark}

\begin{remark}
    The natural pairing
    \[
        V^*\times V \to F,\qquad (\varphi,v)\mapsto \varphi(v),
    \]
    is \(L\)-invariant in the sense that
    \[
        (x\cdot \varphi)(v)+\varphi(x\cdot v)=0
    \]
    for all \(x\in L\), \(\varphi\in V^*\), and \(v\in V\).
\end{remark}

